% ============================================================
% Apêndices (incluídos depois de \appendix)
% ============================================================

\section{Registro de decisões}\label{sec:decisoes}

Cada decisão de método entra aqui com o contexto em que foi tomada, as alternativas
consideradas, a escolha, o motivo e a consequência para o experimento. O registro começa com as
decisões do planejamento e cresce a cada etapa. As decisões mais finas (limiares, listas de
categorias) ficam nos arquivos de configuração e nas seções correspondentes.

{\footnotesize
\setlength{\tabcolsep}{3.5pt}
\renewcommand{\arraystretch}{1.15}
\begin{longtable}{@{}>{\raggedright\arraybackslash}p{0.035\linewidth}
    >{\raggedright\arraybackslash}p{0.155\linewidth}
    >{\raggedright\arraybackslash}p{0.155\linewidth}
    >{\raggedright\arraybackslash}p{0.17\linewidth}
    >{\raggedright\arraybackslash}p{0.21\linewidth}
    >{\raggedright\arraybackslash}p{0.185\linewidth}@{}}
\caption{Registro de decisões.}\label{tab:decisoes}\\
\toprule
\textbf{\#} & \textbf{Contexto} & \textbf{Alternativas} & \textbf{Escolha} & \textbf{Motivo} &
\textbf{Consequência} \\
\midrule
\endfirsthead
\multicolumn{6}{@{}l}{\emph{Tabela~\ref{tab:decisoes}, continuação}}\\
\toprule
\textbf{\#} & \textbf{Contexto} & \textbf{Alternativas} & \textbf{Escolha} & \textbf{Motivo} &
\textbf{Consequência} \\
\midrule
\endhead
\midrule
\multicolumn{6}{r@{}}{\emph{continua}}\\
\endfoot
\bottomrule
\endlastfoot
D1 & Modelo a analisar, entre 4 e 8 bilhões de parâmetros e cabendo numa GPU de 8 GiB. &
Qwen3-4B pós-treinado; Tucano-2b4; GPT-2 small português (piloto). &
Qwen3-4B-Base, revisão fixada por SHA. &
É o modelo de linguagem puro (sem ajuste para diálogo), com arquitetura atual e 36 blocos para a
curva camada a camada. O Tucano tokeniza melhor, mas é menos representativo; o GPT-2 é pequeno e
foi treinado no mesmo texto. &
Português fragmentado (28\% de palavras inteiras): alvos restritos a um token e medidas sempre
por categoria (Seções~\ref{sec:mod-escolha} e~\ref{sec:prob-fragmentacao}). \\
\midrule
D2 & Amostrar cerca de 15 mil vértices que sirvam a P1--P4. &
Parágrafos inteiros (amostragem natural); natural começando pelos alvos. &
Desenho híbrido: núcleo de parágrafos inteiros e estratos de alvos, controles e multitema. &
Parágrafos inteiros deixam só 4 a 9 palavras de conteúdo repetidas em 3 temas e nenhum alvo
utilizável; o híbrido dá cerca de 170 (Seção~\ref{sec:viabilidade}). &
Frequências na amostra refletem o desenho (faixas também pelo corpus); P3 só no núcleo; cerca
de 330 mil tokens a processar. \\
\midrule
D3 & Que vértices tem a rede lexical de tipos. &
Só os $\approx$3,3 mil tipos da amostra; o vocabulário inteiro. &
Vocabulário inteiro (151.643 tipos, sem especiais nem enchimento). &
Descreve a geometria do espaço de entrada sem depender da amostra, e é uma rede grande com
estrutura própria para as métricas da disciplina. &
Uma tabela auxiliar de tipos da amostra continua necessária (não é um recorte;
Seção~\ref{sec:vocab-vs-amostra}); distâncias exatas só na configuração principal. \\
\midrule
D4 & Precisão do cosseno com tolerância de empate $\eps=10^{-6}$ e $d=2560$. &
\emph{float32} na GPU; \emph{float64} na CPU. &
\emph{float64} na CPU, em blocos. &
Em \emph{float32} o erro é de $10^{-6}$ a $10^{-5}$, comparável a $\eps$; empates falsos e
desfeitos mudariam as vizinhanças. &
Cálculo mais lento e cerca de 3 GB de memória na rede do vocabulário; empates exatos e
reprodutíveis. \\
\midrule
D5 & Qual versão do bloco 36 usar, dado que \cod{hidden\_\allowbreak states[36]} é
normalizado no transformers 5.16. &
\cod{hidden\_\allowbreak states[36]} (pós-RMSNorm); saída bruta do bloco. &
Saída bruta por \emph{hook}; a normalizada como robustez (\cod{L36n}). &
Os blocos 1 e 18 só existem na forma bruta; comparar objetos do mesmo tipo na transição
$18\to36$. &
Verificação \cod{norm(hook\_36) == hidden\_states[36]} no modelo real; uma variante a mais na
robustez (Seção~\ref{sec:prob-hidden}). \\
\midrule
D6 & Desempate dos vizinhos na fronteira do $k$-ésimo. &
Só posição (piloto); sorteio; todos os empatados; tipos distintos. &
Sorteio com 20 sementes como principal; as outras três como referência e robustez. &
A posição cria \emph{hubs} artificiais a 15 mil vértices (todas as ocorrências de um tipo
escolhem os mesmos vizinhos). &
Métricas estruturais da rede lexical com média e desvio entre sementes (semente 0 nas
contextuais, onde os empates são raros); piso de ruído e $J^w$ para separar o efeito do
sorteio (Seção~\ref{sec:red-empates}). \\
\midrule
D7 & Biblioteca de redes e detecção de comunidades. &
NetworkX (piloto); igraph; Louvain; Leiden com CPM. &
igraph, com Leiden maximizando a modularidade, 10 sementes; Louvain como checagem. &
Distâncias exatas em C (segundos por rede, contra 20 a 25 minutos no NetworkX); o Leiden
garante comunidades conexas; o padrão CPM do igraph tem outra resolução, por isso o código
passa \cod{objective\_function=\allowbreak"modularity"}. &
NetworkX fica como referência nos testes; estabilidade das comunidades reportada. \\
\midrule
D8 & O que fazer com o piloto e com o código contrafactual de gênero herdado. &
Estender o piloto; reescrever tudo; congelar. &
Piloto congelado (rodando e com testes); código contrafactual e notebook removidos (ficam no
histórico do \cod{git}). &
O piloto é monolítico e não escala, mas documenta os achados que motivaram o método; o código
contrafactual não serve às perguntas atuais. &
Pipeline novo em módulos próprios; um teste de regressão exige que o k-NN novo (posição,
união) reproduza o do piloto. \\
\midrule
D9 & Como configurar execuções (principal e ensaio). &
Argumentos de linha de comando; vários YAMLs por etapa; um YAML por execução. &
Um YAML por execução (\cod{configs/\allowbreak experiment.yaml}; \cod{configs/\allowbreak mini.yaml} herda dele). &
Uma execução fica descrita por um arquivo só, gravado no \cod{\_manifest.json} de cada etapa. &
Artefatos separados por execução (\cod{outputs/\allowbreak experiment/\allowbreak <nome>/}); o ensaio não sobrescreve a
execução principal. \\
\midrule
D10 & Como documentar o experimento. &
Escrever o relatório no fim; relatório vivo. &
Relatório técnico vivo, com figuras, tabelas e números gerados pelo pipeline. &
O texto nunca diverge da última execução, e as decisões e problemas são registrados enquanto
ainda estão frescos. &
Regra de trabalho: uma etapa só está pronta com a seção do relatório atualizada e o documento
compilando. \\
\end{longtable}
}

% ------------------------------------------------------------
\section{Reprodutibilidade}\label{sec:reprodutibilidade}

\paragraph{Ambiente.} Python 3.12 com dependências fixadas em \cod{uv.lock}; instalação com
\cod{uv sync --dev}. As versões efetivamente usadas em cada etapa (Python, numpy, torch, transformers,
igraph, networkx, pandas) são gravadas no \cod{\_manifest.json} da etapa e
reunidas na Tabela~\ref{tab:versoes}. A versão do \cod{mwparserfromhell}, que afeta a
limpeza (Seção~\ref{sec:prob-ref}), ainda não é gravada nos manifestos; até lá, ela é a fixada
em \cod{uv.lock} (0.7.2). O relatório compilado indica a data e a revisão do
\cod{git} da execução que gerou os números (\dataExecucao, \revisaoGit).

\paragraph{Modelo.} \cod{Qwen/Qwen3-4B-Base}, revisão
\cod{\seqsplit{906bfd4b4dc7f14ee4320094d8b41684abff8539}} do Hugging Face (código do próprio transformers, sem código remoto).
Pesos em bf16, \cod{device\_map="auto"} e \cod{max\_memory} de 6 GiB na GPU e 10 GiB na CPU
(\emph{offload} de parte dos blocos). O \emph{offload} muda o tempo, não as contas.

\paragraph{Dados.} A lista de artigos, com \cod{pageid} e \cod{revid}, fica versionada em
\cod{data/corpus/manifest.csv}: o mesmo texto pode ser baixado de novo pelas revisões, mesmo
que os artigos tenham mudado. As respostas cruas da API ficam em cache em \cod{data/raw/wiki}.

\paragraph{Sementes.} Todas vêm do YAML da execução:
\begin{itemize}
    \item \cod{corpus.seed} $=438$: ordem de download dos artigos de cada tema;
    \item \cod{sample.seed} $=438$: um único gerador sorteia, em ordem, o núcleo, o limite por
    tipo, alvos, controles e multitema;
    \item \cod{networks.base\_seed} $=438$: a semente $r$ do desempate da representação
    \cod{rep} com $k$ vizinhos usa \cod{numpy.random.default\_rng([438, crc32(rep), k, r])},
    $r=0,\ldots,19$;
    \item Leiden: 10 execuções com sementes fixas do gerador \cod{random} do Python (usado pelo
    igraph); permutações da linha de base com gerador semeado.
\end{itemize}

\paragraph{Comandos.} Cada etapa lê os artefatos da anterior e pode ser refeita sozinha; uma
etapa cujas saídas já existem é pulada, a menos que se passe \cod{--force}.

\begin{center}
\small
\begin{tabularx}{\linewidth}{@{}>{\ttfamily\raggedright\arraybackslash}p{0.47\linewidth} X@{}}
\toprule
\normalfont\textbf{Comando} & \textbf{O que faz} \\
\midrule
uv run gender-networks corpus & baixa e limpa os artigos (\cod{data/corpus/}) \\
uv run gender-networks sample & monta a amostra híbrida (\cod{sample/}) \\
uv run gender-networks extract -{}-verify & checa as premissas no modelo real (8 sequências) \\
uv run gender-networks extract & extrai as representações (\cod{reps/}) \\
uv run gender-networks knn & candidatos e vizinhos k-NN de todas as variantes (\cod{knn/}) \\
uv run gender-networks metrics & métricas estruturais e comunidades (\cod{metrics/}) \\
uv run gender-networks analyze & análises de P1--P4 (\cod{analysis/}) \\
uv run gender-networks lens & extensão ``vizinhos no vocabulário'' (\cod{lens/}) \\
uv run gender-networks report & figuras, tabelas e \cod{numeros.tex} deste relatório \\
uv run gender-networks all & todas as etapas em ordem (sem a extensão) \\
\bottomrule
\end{tabularx}
\end{center}
Todos aceitam \cod{--config} (padrão \cod{configs/experiment.yaml}); o ensaio de ponta a ponta
usa \cod{--config configs/mini.yaml} (2 temas, cerca de 1.500 vértices), cujos artefatos e
relatório gerado ficam em \cod{outputs/experiment/mini/}. O relatório é compilado com
\cod{latexmk -pdf relatorio.tex} em \cod{report/relatorio/}. O piloto congelado roda com
\cod{uv run python scripts/run\_network\_pilot.py}.

\tabelaopcional[etapa \etapa{report} a partir dos \cod{\_manifest.json} de todas as etapas]
    {versoes}
    {Versões do software, revisão do modelo e tempos de cada etapa na última execução.}
    {tab:versoes}
